% Pursuing Stacks (cotes 134-2 to 134-8): archive cote 134-5, site page 19 (PDF page 20), typescript page 206.
% typescript. Draft compared: ps4.tex:852-914.
% AI-transcribed by Claude Opus 5.5 (claude-opus-5-5, adaptive thinking, effort high), Message Batches API, batch ps2 (msgbatch_01NVwr1Yuo8qJf5rykeGBXPK), 2026-11-12; not yet checked by eye.
\scanpage{134-5/19}
\mspage{206}
settled to get\inserted[typed]{the Quillen factorizations in} the general case,  but this should not be too serious a
difficulty I feel.\over{x}{\unsure{\&}} What\inserted{however} makes me still feel a little unhappy in all this\hand{,}
is\inserted{rather} that I did not get a direct proof for an elementary modelizer being a
closed model category  -  I finally have to make a reduction to the known
case of semisimplicial complexes,  settled by Quillen in his notes.  This
\inserted{detour} looks rather artificial  -  it is the \xout{xxxx}\inserted[typed]{first} instance,  and presumably the
last one,  where the theory I am digging out seems to depend on semi-simpli
-cial techniques,\inserted{techniques} which moreover I dont\inserted{really} know and am not really eager to
swallow.  It's just a prejudice maybe,  \xout{\unsure{butxI}} a block maybe against the
semi-simplicial approach which I never really liked nor assimilated  -
but I do have the feeling that the more refined and specific semi-simplici-%
al techniques and notions (such as minimal fibrations,  used in Quillen's
proof though,  alas !)  are \xout{\illegible{1}} \hand{irrelevan}t for an understanding of
the main structures entering homotopy theory and homo\hand{to}pical algebra.  As
for the notion of a Kan complex or a Kan fibration  -  namely just a
``fibration'' in Quillen's axiomatic set-up,which I was \deleted{\xout{\illegible{1}}} finally glad to
find\hand{,}``ready foruse''  -  I \deleted{finally} came to convince myself\inserted{at last} that it was a
basic notion indeed,  and it was no use trying to by\deleted{e}-pass it at all price.
Thus I took t\over{\unsure{b}}{o} the opposite,  and tried to \xout{\unsure{develop}} pin down a Quillen\inserted{-}type
factorization theorem,  and \xout{the} his\hand{cha}racteristic sea-saw game \xout{of}\inserted[typed]{between} right
and left lifting properties,  in \deleted{the}\inserted{as} great\deleted{est} generality\inserted{as} I could manage.

30.5.

71.  The \deleted{bulk of} scratchwork done since last month has of course conside-%
rably cleared up the prospects of my present pre-stacks reflection on
homotopy models,  in which I unsuspectingly embarked three months ago.
\over{\illegible{1}}{Now} I would like to sketch a provisional working program for the notes
still ahead.

A)  Write down at last the story of aspericity structures and canonical
modelizers,  as I was about to when I interrupted the notes to do my \deleted{scrat}
scratchwork.

B)  Study the basic modelizer (Cat),  and the common properties of
elementary modelizers $A^{\wedge}$,  with a main emphasis upon base change and
co-base change properties,  and upon Quillen type factorization questions.
Here it will be useful to dwell somewhat on the ``homotopy integral'' variant
of taking amalgamated sums in (Cat),  on the analoguous constructions for
topoi,  and how these compare to the usual amalgamated sums, including the
interesting case of topological spaces.  It turns out that the homotopy
\margin{19}
